\begingroup
\small
\setlength{\parindent}{0pt}
\setlength{\parskip}{0.55\baselineskip}
\setlength{\abovedisplayskip}{8pt plus 2pt minus 2pt}
\setlength{\belowdisplayskip}{8pt plus 2pt minus 2pt}
\setlength{\jot}{4pt}

\section*{How the Proof Works}

The proof can be seen in one continuous calculation.  Set
\[
V_r:=\partial_t^r u,
\qquad
Q_r:=\partial_t^r p,
\qquad
\Lambda:=(-\Delta)^{1/2}.
\]
Differentiating Navier--Stokes \(r\) times gives
\[
\partial_tV_r
+\sum_{a=0}^{r}\binom ra(V_a\!\cdot\nabla)V_{r-a}
+\nabla Q_r
=\nu\Delta V_r.
\]
At spatial Sobolev height \(s\), define
\[
E_{r,s}(t):=\frac12\|\Lambda^sV_r(t)\|_2^2.
\]
Apply \(\Lambda^s\) to the differentiated equation and pair it with
\(\Lambda^sV_r\).  The time term and the viscous term are
\[
\left\langle\Lambda^s\partial_tV_r,\Lambda^sV_r\right\rangle
=\partial_tE_{r,s},
\qquad
\nu\left\langle\Lambda^s\Delta V_r,\Lambda^sV_r\right\rangle
=-2\nu E_{r,s+1}.
\]
Keep the nonlinear and pressure terms together as
\[
\mathcal P_{r,s}
:=-\operatorname{Re}
\left\langle
\Lambda^sV_r,
\Lambda^s\!\left[
\sum_{a=0}^{r}\binom ra(V_a\!\cdot\nabla)V_{r-a}
+\nabla Q_r
\right]
\right\rangle.
\]
The pairing gives the exact identity
\[
\boxed{
\partial_tE_{r,s}=\mathcal P_{r,s}-2\nu E_{r,s+1},
\qquad
E_{r,s+1}
=\frac{\mathcal P_{r,s}-\partial_tE_{r,s}}{2\nu}.
}
\]
This is the central move.  The time variation of the energy at height \(s\)
contains the energy at height \(s+1\), and viscosity supplies the coefficient
that connects them.  Differentiating again repeats the same move; after \(k\)
steps,
\[
\boxed{
\partial_t^kE_{r,s}
=(-2\nu)^kE_{r,s+k}
+\sum_{j=0}^{k-1}(-2\nu)^j
\partial_t^{\,k-1-j}\mathcal P_{r,s+j}.
}
\]
At the velocity level the same coupling is visible directly:
\[
V_{r+1}
=\nu\Delta V_r
-\sum_{a=0}^{r}\binom ra(V_a\!\cdot\nabla)V_{r-a}
-\nabla Q_r,
\qquad
\nu\Delta:H^{m+1}\longrightarrow H^{m-1}.
\]
One further time derivative thus carries the parabolic weight of two spatial
derivatives, while one time derivative of the quadratic Sobolev energy raises
its height by one.  This is why the estimates form a two-dimensional array:
temporal order labels the rows, spatial Sobolev height labels the columns, and
viscosity links neighboring positions.

Now apply the mechanism at any fixed finite depth \(N\) and height \(M\).
Theorem~\ref{thm:datum-rectangles} carries the required adjacent rows across
the whole interval \((0,T_*)\):
\[
V_n\in L^2(0,T_*;H^{M+1}),
\qquad
V_{n+1}\in L^2(0,T_*;H^{M-1})
\qquad(0\le n\le N).
\]
Finite summation and Proposition~\ref{prop:datum-terminal-realization} give
\[
\mathfrak R^*_{N,M}[u_0,\nu,T_*]
=\sum_{n=0}^{N}
\left(
\|V_n\|_{L^2_tH^{M+1}_x}^2
+\|V_{n+1}\|_{L^2_tH^{M-1}_x}^2
\right)<\infty.
\]
The pressure recurrence supplies the matching \(Q_n\)-coordinates, so this
finite calculation is the rectangle \(\mathcal R_{N,M}[u_0;T_*]\).  Enlarging
\((N,M)\) leaves every shared coordinate unchanged.  The rectangles can thus
be taken together, and for every finite \(r,m\),
\[
\mathcal I[u_0]
\in\varprojlim_{N,M}\mathscr X_{N,M}
\quad\Longrightarrow\quad
u\in H^r(0,T_*;H^m(\mathbb R^3)).
\]

\clearpage

Now ask for any downstream estimate whose hypotheses mention finitely many
derivatives of \(u\).  Write those derivatives as
\[
D_j u:=\partial_t^{a_j}\nabla^{b_j}u,
\qquad 1\le j\le J,
\]
where \(J\), every \(a_j\), and every \(b_j\) are finite.  Suppose the desired
norm for \(D_j u\) is \(L_t^{p_j}L_x^{q_j}\).  Choose finite Sobolev indices
\(n_j,m_j\) so that
\[
n_j-a_j>\frac12-\frac1{p_j},
\qquad
m_j-b_j>3\!\left(\frac12-\frac1{q_j}\right).
\]
The time and space Sobolev embeddings then give
\[
H^{n_j}(0,T_*;H^{m_j}(\mathbb R^3))
\hookrightarrow
W^{a_j,p_j}(0,T_*;W^{b_j,q_j}(\mathbb R^3)),
\]
and hence
\[
\|D_j u\|_{L_t^{p_j}L_x^{q_j}}
\le
C_j\|u\|_{H_t^{n_j}H_x^{m_j}}<\infty.
\]
Because the requested list is finite, one can choose
\[
N\ge\max_{1\le j\le J}n_j,
\qquad
M\ge\max_{1\le j\le J}m_j,
\]
and a single rectangle \(\mathcal R_{N,M}[u_0;T_*]\) contains every coordinate
used by the estimate.  Nothing is assembled from unrelated bounds: the
coordinates already agree inside the compatible family.

If the requested expression contains products, choose \(M\) above the
Sobolev algebra threshold and apply the standard product or Moser estimate.
If it contains pressure, use
\[
Q_r
=R_iR_j\sum_{a=0}^{r}\binom ra V_{a,i}V_{r-a,j},
\]
together with the same product estimate and the boundedness of the Riesz
transforms.  If it asks for an intermediate exponent, interpolate between two
coordinates already present in the rectangle.  If it asks for time
integrability on \((0,T_*)\), use the time Sobolev embedding or finite-interval
H\"older after selecting the required temporal order.

Thus every finite-order Sobolev-controlled request has the same form:
\[
\boxed{
\begin{gathered}
\text{finite derivative and norm requirements}
\longrightarrow
\text{one sufficiently large }(N,M)\\
\longrightarrow
\mathcal R_{N,M}[u_0;T_*]
\longrightarrow
\text{the requested estimate}.
\end{gathered}
}
\]
In precise terms, if a quantity \(\mathfrak E[u;T_*]\) is bounded by finitely
many such Sobolev coordinates,
\[
\mathfrak E[u;T_*]
\le
F\!\left(
\|u\|_{H_t^{n_1}H_x^{m_1}},\ldots,
\|u\|_{H_t^{n_J}H_x^{m_J}}
\right)
\]
for a finite-valued continuous function \(F\), then
\[
\boxed{
\mathcal I[u_0]
\longrightarrow
\mathcal R_{N,M}[u_0;T_*]
\longrightarrow
\mathfrak E[u;T_*]<\infty.
}
\]
That is the general rule.  The estimate may be a pointwise derivative bound,
a spacetime norm, a nonlinear product bound, a pressure bound, or a finite
collection of them.  The chosen \((N,M)\) and its constant may change with the
question.  No estimate uniform over all derivative orders is being claimed or
needed.

The endpoint is obtained by using the same rule at each fixed height.  For
every finite \(m\), the intersection supplies enough spatial control for the
equation to give
\[
u\in W^{1,1}(0,T_*;H^m)
\hookrightarrow C([0,T_*];H^m).
\]
The endpoint values agree under the Sobolev inclusions because they are all
limits of the same function \(u(t)\).  Consequently,
\[
u_*:=u(T_*)\in\bigcap_{m<\infty}H^m=H^\infty.
\]
Local well-posedness continues the solution from \(u_*\) past \(T_*\), which
contradicts finite maximality and gives \(T_*=\infty\).

The meaning is simple.  The intersection is not one enormous norm with one
constant.  It says that every finite regularity question can be answered
inside a finite rectangle, and that all such answers concern the same
solution and agree wherever they overlap.  Whatever finite classical estimate
is requested, first count the derivatives and integrability it requires, then
choose a rectangle large enough to contain them.  The estimate is extracted
from that rectangle by ordinary functional analysis.  This is why the
estimates come from the intersection rather than producing it.

\endgroup
